\subsection{Range and Domain of an Absolute Value Function}
Absolute value functions have a simple domain.
\begin{fact}[Domain of Absolute Value Functions]
The domain of $f(x)=\Abs{x}$ is all real numbers (we can take the absolute value of any real number). Composing with linear functions does not change the domain, so the domain of any absolute value function is all real numbers.
\end{fact}
The range is more challenging.
\begin{fact}[Range of the Base Absolute Value]
The absolute value is always \makeblank[1in] or \makeblank[1in]. So if $f(x)=\Abs{x}$, $$\rng f=\makeblank$$
\end{fact}
For the range of a general absolute value function, it is easiest to consider the graph.
\begin{example}Graph $f(x)=\frac{1}{3}\Abs{x+1}-1$. Use this graph to determine the range.
\begin{lpic}{}{8}
\begin{scope}[xshift=\value{cols}*2\linespace-6\linespace,yshift=-4\linespace]
\setgraph{5}{5}{3}{3}
\setspace{1}
\AxesLarge
\end{scope}
\regtext{1.5}{-4}{$x$};
\regtext{5.5}{-4}{$y=f(x)$};
\draw (0,-4) -- (8,-4);
\draw (3,-3) -- (3,-7);
\regtextleft{1}{-8}{$\rng f={}$};
\end{lpic}
\end{example}
\begin{fact}[Range of an Absolute Value Function]
For an absolute value fuction $f(x)=a\Abs{bx+c}+d$:
\begin{itemize}
\item If $a>0$, the ``V'' shape opens \makeblanksmall, and $\rng f=\makeblank[1in]$
\item If $a<0$, the ``V'' shape opens \makeblanksmall, and $\rng f=\makeblank[1in]$ 
\end{itemize}
\end{fact}